\documentclass[11pt,a4paper]{article}
\usepackage{gtu-el-yazisi}
\begin{document}
\notbilgi{Sayısal Haberleşme}{Matched filter, SAK ve 1. Nyquist kriteri}{Ders-10.pdf}

\sayfa{1}
\textcolor{kirmizikalem}{genel}
\begin{gather*}
\gamma_0=\frac{\sigma_0^2}{a_1-a_2}\,\ln\frac{P(s_2)}{P(s_1)}+\frac{a_1+a_2}{2}\;\rightarrow\\
P_b=\Bigl[1-Q\Bigl(\frac{\gamma_0-a_1}{\sigma_0}\Bigr)\Bigr]P(s_1)+Q\Bigl(\frac{\gamma_0-a_2}{\sigma_0}\Bigr)P(s_2)
\end{gather*}
\[
\Big\downarrow\qquad P(s_1)=P(s_2)=\frac{1}{2}
\]
\[
\gamma_0=\frac{a_1+a_2}{2}\;\rightarrow\;
P_b=Q\Bigl(\frac{a_1-a_2}{2\sigma_0}\Bigr)=Q\Biggl(\sqrt{\frac{E_d}{2N_0}}\Biggr)=Q\Biggl(\sqrt{\frac{E_b}{N_0}(1-\rho)}\Biggr)
\]
\[
a_i=\int_0^T s_i(t)\bigl[s_1(t)-s_2(t)\bigr]dt=\underbrace{s_i(t)\circledast h(t)}_{a_i(t)}\Big|_{t=T}
\]
\elcizimi[0.85]{sekiller/Ders-10-s1-1.png}[1]
\[
h(t)=s_1(T-t)-s_2(T-t)
\]
\textcolor{kirmizikalem}{! Formüllerin nasıl çıktığını öğren}

Soruda $h(t)$ $s_1(t)$'yi karşılıyor dersek \quad $h(t)=s_1(T-t)$ olur integral devresinde
\elcizimi[0.2]{sekiller/Ders-10-s1-2.png}[1]
\begin{align*}
\sigma_0^2&=\frac{N_0}{2}E_h\\
E_d&=\int_0^T\bigl[s_1(t)-s_2(t)\bigr]^2dt\qquad\qquad E_b=\sum_{i=1}^{2}P(s_i)\,E_{s_i}\\
\rho&=\frac{1}{E_b}\int_0^T s_1(t)\,s_2(t)\,dt
\end{align*}
\begin{ornek}[(Unipolar signaling)]
``1'' $\rightarrow s_1(t)=A,\ 0\le t\le T$\\
``0'' $\rightarrow s_2(t)=0,\ 0\le t\le T$\\
$P(s_1)=\dfrac{1}{4},\ P(s_2)=\dfrac{3}{4}$\\
$A=1\ \mathrm{V},\ T=10^{-3}\ \mathrm{sn},\ \dfrac{N_0}{2}=10^{-3}\,\dfrac{\mathrm{Watt}}{\mathrm{Hz}}$\\
$\gamma_0=?\quad P_b=?$
\end{ornek}
\begin{align*}
a_1&=A^2T=1^2\cdot 10^{-3}=10^{-3}\\
a_2&=0\\
h(t)&=s_1(T-t)\underbrace{-\,s_2(T-t)}_{0}=s_1(T-t)
\end{align*}
\elcizimi[0.95]{sekiller/Ders-10-s1-3.png}[1]

\sayfa{2}
\begin{align*}
E_h&=\int_0^T h^2(t)\,dt=\int_0^T A^2\,dt=A^2T=10^{-3}\\
\sigma_0^2&=\frac{N_0}{2}\cdot E_h=10^{-3}\cdot10^{-3}=10^{-6}\\
\gamma_0&=\frac{\sigma_0^2}{a_1-a_2}\ln\frac{P(s_2)}{P(s_1)}+\frac{a_1+a_2}{2}\\
&=\frac{10^{-6}}{10^{-3}}\ln\frac{3/4}{1/4}+\frac{10^{-3}+0}{2}=1{,}5986\cdot10^{-3}\simeq1{,}6\cdot10^{-3}
\end{align*}
\begin{multline*}
P_b=\Bigl[1-Q\Bigl(\frac{\overset{\textcolor{turuncukalem}{1{,}6\cdot10^{-3}}}{\gamma_0}-\overset{\textcolor{turuncukalem}{10^{-3}}}{a_1}}{\underset{\textcolor{turuncukalem}{10^{-3}}}{\sigma_0}}\Bigr)\Bigr]\underset{\textcolor{turuncukalem}{\frac14}}{P(s_1)}\\
+Q\Bigl(\frac{\overset{\textcolor{turuncukalem}{1{,}6\cdot10^{-3}}}{\gamma_0}-\overset{\textcolor{turuncukalem}{0}}{a_2}}{\underset{\textcolor{turuncukalem}{10^{-3}}}{\sigma_0}}\Bigr)\underset{\textcolor{turuncukalem}{\frac34}}{P(s_2)}
\end{multline*}
\begin{gather*}
P_b=\Bigl[1-\underbrace{Q(0{,}6)}_{2{,}7426\cdot10^{-1}}\Bigr]\frac14+\underbrace{Q(1{,}6)}_{5{,}4799\cdot10^{-2}}\cdot\frac34\\
=0{,}2225
\end{gather*}
\begin{ornek}
Transfer function of a MF is $H(f)=\dfrac{2-2e^{-j2\pi fT}}{j2\pi f}$\\
a) Find impulse response of MF and sketch it.\\
b) $''$ the pulse $s(t)$ matched by this MF an sketch it.\\
c) If this $s(t)$ is applied to the input of MF, find the output signal of the MF and sketch it. Find the value at the sample instant.\\
d) If this $s(t)$ is applied to the input of the corelator, find the output signal of the correlator and sketch it. Find the value at the simple instant.
\end{ornek}
\[
\text{a)}\ H(f)=\frac{2}{j2\pi f}-\frac{2e^{-j2\pi fT}}{j2\pi f}\ \xrightarrow{\ \mathcal{F}^{-1}\ }\ h(t)=\mathrm{sgn}(t)-\mathrm{sgn}(t-T)
\]
\elcizimi[0.55]{sekiller/Ders-10-s2-1.png}[2]
\elcizimi[0.4]{sekiller/Ders-10-s2-2.png}[2]
\[
\text{b)}\quad h(t)=s(\underbrace{T-t}_{t'})\qquad\rightarrow\qquad h(T-t')=s(t')
\]
\[
\begin{array}{l@{\qquad\qquad}l}
t'=T-t & \Longrightarrow\ \kutula{s(t)=h(T-t)}\\
t=T-t' &
\end{array}
\]
\elcizimi[0.4]{sekiller/Ders-10-s2-3.png}[2]

\sayfa{3}
\elcizimi[0.4]{sekiller/Ders-10-s3-1.png}[3]
\elcizimi[0.65]{sekiller/Ders-10-s3-2.png}[3]
\[
z_1(t)=s_1(t)\circledast h(t)=s(t)\circledast s(t)
\]
\[
A\,\Pi\!\Bigl(\frac{t}{T}\Bigr)\circledast A\,\Pi\!\Bigl(\frac{t}{T}\Bigr)=A^2T\,\Lambda\!\Bigl(\frac{t}{T}\Bigr)
\]
\elcizimi[0.35]{sekiller/Ders-10-s3-3.png}[3]
\elcizimi[0.5]{sekiller/Ders-10-s3-4.png}[3]
\[
z_1(T)=4T
\]
\elcizimi[0.35]{sekiller/Ders-10-s3-5.png}[3]
\elcizimi[0.35]{sekiller/Ders-10-s3-6.png}[3]
\elcizimi[0.35]{sekiller/Ders-10-s3-7.png}[3]
\[
z_1(t)=s(t)\circledast h(t)=\int_{-\infty}^{\infty}s(z)\,h(t-z)\,dz
\]
\begin{align*}
t<0&:\quad z_1(t)=0\\
0<t<T&:\quad z_1(t)=\int_0^t 2\cdot 2\,dt=4t
\end{align*}

\sayfa{4}
\elcizimi[0.4]{sekiller/Ders-10-s4-1.png}[4]
\begin{align*}
T<t<2T&:\quad z_1(t)=\int_{t-T}^{T}2\cdot 2\,dt=\\
t>2T&:\quad z_1(t)=0
\end{align*}
\[
z_1(t)=\begin{cases}
0, & t<0\\
4t, & 0<t<T\\
4(2T-t), & T<t<2T\\
0, & t>2T
\end{cases}
\]
\elcizimi[0.5]{sekiller/Ders-10-s4-2.png}[4]
\[
z_1(T)=4T
\]
\elcizimi[0.7]{sekiller/Ders-10-s4-3.png}[4]
\begin{align*}
z_2(t)&=\int_0^t s(z)\,s(z)\,dz\\
&=\int_0^t s^2(z)\,dz=4t\qquad\qquad t<T\\
&\phantom{=}\int_0^T s^2(z)\,dz=4T\qquad\qquad t>T
\end{align*}
\[
z_2(t)=\begin{cases}
4t, & t<T\\
4T, & t>T
\end{cases}
\]
\elcizimi[0.45]{sekiller/Ders-10-s4-4.png}[4]
\[
z_2(T)=4T
\]
\section{Semboller arası karışım (SAK)}
(Inter symbol Interference - ISI)

(Pulse shaping - Darbe şekillendirme)

\elcizimi[0.95]{sekiller/Ders-10-s4-5.png}[4]

\sayfa{5}
\elcizimi[0.95]{sekiller/Ders-10-s5-1.png}[5]
\[
H(f)=H_t(f)\,H_c(f)\,H_r(f)
\]
\section{Sıfır SAK için 1. Nyquist kriteri}
(1st Nyquist criterion for zero ISI)
\[
h(t)=\begin{cases}1, & t=0\\ 0, & t=\pm nT\end{cases}
\;\Rightarrow\;
h(nT)=\begin{cases}1 & n=0\\ 0 & n\neq0\end{cases}
\]
\[
\begin{array}{l}
n=1,2,3,\ldots\\
T=\dfrac{1}{R_s}\\
\hspace{1.5em}\uparrow\ \text{sembol hızı}
\end{array}
\qquad\qquad
\begin{array}{l}
\text{Hız}=R_s\,\dfrac{\text{sembol}}{\text{sn}}\\
\Rightarrow\ BG=\dfrac{R_s}{2}\ \text{Hz}
\end{array}
\]
\[
\xrightarrow{\qquad\mathcal{F}^{-1}\qquad}\quad h(t)=\mathrm{sinc}\Bigl(\frac{t}{T}\Bigr)
\]
\elcizimi[0.55]{sekiller/Ders-10-s5-2.png}[5]
\[
A\tau\,\mathrm{sinc}(t\tau)\ \overset{\mathcal{F}}{\longleftrightarrow}\ A\,\Pi\!\Bigl(\frac{f}{\tau}\Bigr)
\]
\textcolor{kirmizikalem}{$\leftarrow$ ideal Nyquist filter for zero ISI}

\textcolor{murekkep}{mavinin katkısı 1, siyahın 0} \quad semboller arası karışım 0
\elcizimi[0.9]{sekiller/Ders-10-s5-3.png}[5]
\textcolor{turuncukalem}{Senkron hatasını yaptığımızı düşünelim sembol arası karışım olabilir.}
\[
\frac{t}{T}=k\in\mathbb{Z}-\{0\}\;\Rightarrow\;t=kT
\]
\section{Raised cosine (RC) darbesi}
Daha yüksek azalma hızına sahiptir \quad senkronizasyon hatası azalır
\[
BG=k\,\frac{R_s}{2}\qquad 1\le k\le 2
\]
\section{1. Nyquist kriterinin diğer ifadesi}
\begin{multline*}
\underset{\substack{\text{sampled}\\ h(t)}}{h_s(t)}=h(t)\,\overset{\text{\okunamadi{darbe} katarı,}}{s_T(t)}=h(t)\sum_{k}\delta(t-kT)\\
=\sum_{k}h(kT)\,\delta(t-kT)=\delta(t)
\end{multline*}

\sayfa{6}
\[
\mathcal{F}\Bigl\{h(t)\sum_{k}\delta(t-kT)\Bigr\}=\mathcal{F}\bigl\{\delta(t)\bigr\}
\]
\[
H(f)\circledast\frac{1}{T}\sum_{k}\delta\Bigl(f-\frac{k}{T}\Bigr)=1
\]
\[
\frac{1}{T}\sum_{k}H\Bigl(f-\frac{k}{T}\Bigr)=1\;\Rightarrow\;\kutula{\sum_{k}H\Bigl(f-\frac{k}{T}\Bigr)=T}
\]

\end{document}
